% Documento creado por �scar Varela Su�rez y Jos� Antonio Almena Mu�oz
% Haz con �l lo que quieras...

% si queremos que un cap�tulo nuevo empiece en cualquier p�gina (dcha. o izda.) ponemos dentro de las
% opciones de documentclass --> openany
\documentclass[twoside,a4paper,12pt]{book}

\usepackage[dvips]{graphicx}
\usepackage[spanish]{babel}
%\selectlanguage{spanish}
%\usepackage[T1]{fontenc}
\usepackage[latin1]{inputenc}
\usepackage{babel}
\usepackage{fancyhdr}
\usepackage{caption2}
\usepackage{geometry}
\usepackage{makeidx}
\usepackage{graphicx}
\usepackage{hthtml}
\usepackage{latexsym}
\usepackage{amssymb}
\usepackage{eucal}
\usepackage{setspace}\singlespacing
% Para otros interlineados: \renewcommand{\baselinestretch}{1.05}
\usepackage{dsfont} %Para las formulitas

%\pagestyle{headings}

\geometry{left=2.5cm, right=2.5cm, top=2.5cm, bottom=2.5cm}

%para crear un �ndice alfab�tico (de usarlo lo hacemos al final del documento)
\makeindex

\begin{document}

\renewcommand{\captionlabelfont}{\textbf}

\pagestyle{fancy}

\fancyhf{}
%
\renewcommand{\headrulewidth}{0.5pt}

\fancyhead[LO]{\rightmark} % En las p�ginas impares, parte izquierda del encabezado, aparecer� el nombre de cap�tulo

\fancyhead[RE]{\leftmark} % En las p�ginas pares, parte derecha del encabezado, aparecer� el nombre de secci�n

\fancyhead[RO,LE]{\thepage} % N�meros de p�gina en las esquinas de los encabezados
%

% para no numerar la p�gina
\thispagestyle{empty}

% si se comenta la siguiente l�nea en el encabezado de una pagina aparece el cap�tulo al que pertenece en lugar
% de la secci�n a la que pertenece
\baselineskip 1.35\baselineskip

\vspace{2cm}
%
\begin{figure}[htb]
%%Escudo de la reyjuan
%Para que inserte la imagen en el pdf ejecutar DVI2DPF
\centerline{\resizebox{.05\textwidth}{!}{\includegraphics{logo_urjc.eps}}}
\end{figure}
%
\begin{center}
{\Large {\bf Universidad Rey Juan Carlos}} \vspace{5mm}

%
{\large Escuela Superior de Ciencias Experimentales y Tecnolog�a}
\vspace{5mm}

%
{\large Departamento de Inform�tica, Estad�stica y Telem�tica}

%
\vspace{4.5cm}
%
%
%

{\Large {\bf ------ t�tulo ------}}

\vspace{3cm}

%
{\large {\bf Trabajo Fin de Carrera}}
%
%
\vspace{4cm}


--- Autor ---\\
\vspace{1cm}
---Tutor---\\
\vspace{1cm} \today

%
\end{center}
%Para que nos deje una p�gina en blanco despu�s de la portada
\newpage{\pagestyle{empty}\cleardoublepage}


\setcounter{page}{1} %para que no cuente la primera p�gina (la portada)
%\setcounter{secnumdepth}{2} \setcounter{tocdepth}{2}


% el documento est� desglosado en diferentes ficheros
% para incluir cada uno de ellos en el total se utiliza la sentencia: \include{nombre_fichero}

% para que no ponga "Cap�tulo 0" en las secciones agradecimientos y resumen
\renewcommand{\chaptermark}[1]{\markboth{\small{\  #1}}{}}

\frontmatter

% ----------------- AGRADECIMIENTOS --------------------
\chapter{Agradecimientos}
%\thispagestyle {empty}

Agradezco toda la ayuda prestada por \dots %\dots pone los puntos suspensivos: "..."

% ----------------- RESUMEN ----------------------------
\chapter{Resumen}
%\thispagestyle {empty}

    Resumen del documento \dots


% para que vuelva a poner "Cap�tulo n" donde corresponda a partir de esta l�nea
\renewcommand{\chaptermark}[1]{\markboth{\small{\chaptername \ \thechapter. #1}}{}}
% ----------------- �NDICE GENERAL ---------------------
\tableofcontents
% ----------------- �NDICE DE FIGURAS ------------------
\listoffigures

\mainmatter

% ----------------- INTRODUCCI�N -----------------------
\chapter{Introducci�n (10\%)}

Aqu� har�s una breve introducci�n del problema a resolver. Para
ello se habla de:

\section {qu� quieres resolver}
\dots
\section {c�mo lo vas a hacer}
\dots
\section {herramientas conceptuales necesarias}
 \dots
\section {herramientas utilizadas}
\dots

% ----------------- ESTADO DEL ARTE --------------------
\chapter{Estado del arte (10-20\%)}
Bas�ndose en art�culos, libros, etc. que se os haya facilitado y
de otros que estim�is oportuno, se hablar� de:
\section{Descripci�n del problema}
\dots
\section{Descripci�n de los trabajos anteriores que se han dedicado a resolverlo}
\dots

% ----------------- OBJETIVO ---------------------------
\chapter {Objetivo (5\%)}

\section {Descripci�n, en un objetivo general, de la finalidad del proyecto.}
\dots
\section {Descripci�n de objetivos parciales que se necesitan cubrir para llegar al objetivo final}
\dots
\section {Descripci�n de alto nivel de las etapas que sigues en el desarrollo}
\dots

% ----------------- DESCRIPCI�N INFORM�TICA ------------
\chapter{Descripci�n Inform�tica (20-35\%)}

Para ello describir�s:

\section{La base de datos coleccionada (si tiene sentido).}
\dots
\section{Los algoritmos para el desarrollo de la soluci�n}
\dots

% ----------------- RESULTADOS EXPERIMENTALES ----------
\chapter{Resultados Experimentales (20-35\%)}

En este apartado deber�n quedar reflejados los experimentos
realizados. Para ello se mostrar�n:

\section {Resultados en forma de tablas, gr�ficas e im�genes donde se describa cuantitativa
            y cualitativamente el funcionamiento de la aplicaci�n} \dots

\section {An�lisis cr�tico de los resultados con el objetivo de decidir si el sistema implementado es v�lido}
\dots

% ----------------- CONCLUSIONES Y TRABAJOS FUTUROS ----
\chapter{Conclusiones y trabajos futuros (5\%)}

Resumen de los logros principales conseguidos, destacando:

\section {Implementaci�n}
\dots

\section {Resultados}
\dots

En futuros trabajos, a partir de una cr�tica constructiva del
trabajo realizado, plantear mejoras y extensiones del mismo.


\backmatter

% para que no ponga "Cap�tulo n" en las secciones anexo y bibliograf�a
\renewcommand{\chaptermark}[1]{\markboth{\small{\  #1}}{}}
% ----------------- ANEXOS -----------------------------
\chapter{Anexo 1}
Lo relativo al anexo 1 \dots

Ejemplos de tratamiento de texto:

Probamos una cita \cite{NewCam97}\\
Citamos un libro\cite{Pesce}\\
Probamos a poner una nota al pi�\footnote{Mi primera nota al pi�}\\
Probando: \it{cursiva} \textbf{negrita} \underline{subrayada} \emph{enfatizar}\\


% pueden hacer falta varios anexos
%\chapter{Anexo 2}

% ----------------- BIBLIOGRAF�A -----------------------
    % para que funcione bien esto se debe compilar all� donde halla una cita usando "bibtext"
    % el "alpha" es para que ordene la bibliograf�a por orden alfab�tico
\bibliographystyle{alpha} \bibliography{librero}
% incluye la bibliograf�a en el �ndice
\addcontentsline{toc}{chapter}{Bibliograf�a}

%\printindex  %%%% crea un �ndice alfab�tico con las palabras indexadas en el documento: \index{palabra}

\end{document}
